\documentclass[]{kyvernitis-resume}
\fullname{Aryana Haghjoo}

\begin{document}
\resumeheader
{\email{aryana.haghjoo@email.ucr.edu}}
{\website{aryana-haghjoo.github.io}}
{\linkedin{aryana-haghjoo}}
{\github{aryana-haghjoo}}
\medskip

Aryana Haghjoo is an early-career astrophysicist with a strong foundation in both theoretical and data-driven research. She holds an M.Sc. in Physics from McGill University, where she specialized in computational global 21cm cosmology, and is currently pursuing a Ph.D. in Astronomy at the University of California, Riverside. Her recent research leverages deep learning to super-resolve low-resolution galaxy spectra.
%She is currently extending this work to a multi-modal foundation model that integrates photometric and spectroscopic data for enhanced scientific inference across large-scale surveys.\par


\medskip
%--------------------------------------------------
\begin{section}{Education}
    \begin{subsectionnobullet}{Ph.D. Astronomy, \href{https://www.ucr.edu}{University of California, Riverside (UCR)}}{GPA: 3.7/4}{Sept 2025 -- Sept 2028 (Expected)}{Riverside, California, USA}
    \end{subsectionnobullet}

    \begin{subsectionnobullet}{M.Sc.\ Astronomy, \href{https://www.ucr.edu}{UCR}}{GPA: 3.7/4}{Sept 2024 -- Jun 2026}{Riverside, California, USA}
    \end{subsectionnobullet}
    
    \begin{subsectionnobullet}{M.Sc.\ Physics, \href{https://www.physics.mcgill.ca/}{McGill University}}{GPA: 3.5/4}{Sept 2021 -- Dec 2023}{Montreal, Quebec, Canada}
        Thesis: \href{https://www.proquest.com/openview/e3346eb2737e36a5837fe4141e6ad92a/1?pq-origsite=gscholar\&cbl=18750\&diss=y}{Parameter Estimation of the Global 21cm Signal}
    \end{subsectionnobullet}
    
    \begin{subsectionnobullet}{B.Sc.\ Physics, \href{http://physics.sharif.ir/~phyweb/}{Sharif University of Technology (SUT)}}{GPA: 18.25/20}{Sept 2018 -- Jul 2021}{Tehran, Iran}
    \end{subsectionnobullet}
\end{section}
%--------------------------------------------------
%--------------------------------------------------
\begin{section}{Publications}

\textit{14 publications, 106 citations, h-index: 6 (NASA ADS, Oct 2026).}

\vspace{0.2cm}

\textbf{First-Author Manuscripts}

\begin{itemize}

\item \textbf{Haghjoo, A.}, Hemmati, S., Mobasher, B., Chartab, N., de la Vega, A., Eifler, T., Ramey, E., Nayyeri, H., \& Sattari, Z.
\textit{Learning to See Sharper: A Physics-Informed Artificial Intelligence Framework for Super-Resolving Galaxy Spectra.}
Under review at the \textit{Astrophysical Journal}; \href{https://arxiv.org/abs/2603.18357}{arXiv:2603.18357}.

\item \textbf{Haghjoo, A.}, Hemmati, S., \& Mobasher, B.
\textit{Sharper at What Cost? Benchmarking Deep Learning Against Classical Deconvolution for Galaxy Spectral Super-Resolution.}
Under review at the \textit{Astrophysical Journal}.

\item \textbf{Haghjoo, A.}, Hemmati, S., Mobasher, B., Chartab, N., \& Sattari, Z.
\textit{Spectral Super-Resolution for Roman Grism Spectroscopy.}
In preparation for submission to the \textit{Astrophysical Journal}.

\item \textbf{Haghjoo, A.}, Mobasher, B., \& D'Aloisio, A.
\textit{Connecting Observational Spectra to the Bubble Size Distribution at the Epoch of Reionization with Deep Learning.}
In preparation for submission to the \textit{Astrophysical Journal}.


% \item \textbf{Haghjoo, A.}, \& Barish, B.
% \textit{The Global 21 cm Signal as a Probe of the Early Universe: A Review of Theory and Observation Attempts.}
% In preparation for submission to the \textit{Astrophysical Journal}.

\end{itemize}

\vspace{0.2cm}

\textbf{Selected Co-Author Publications}

\begin{itemize}

\item Taamoli, S., Mobasher, B., Hatamnia, H., \textbf{Haghjoo, A.}, et al.
\textit{The Role of the Cosmic Web in Galaxy Evolution Across $0.4<z<4$.}
\textit{arXiv preprint} (2026).

\item Petric, A., Noirot, G., Alberts, S., \textbf{Haghjoo, A.}, et al.
\textit{Surveying the Universe in 4D: Beating Cosmic Variance with Wide-Field Slitless Spectroscopy from HST, JWST, Euclid, Roman, and Beyond.}
\textit{arXiv preprint} (2026).

\item Zhang, Y., Ding, X., Yang, L., \textbf{Haghjoo, A.}, et al.
\textit{Extended Components of Little Red Dots in the Rest-Frame Optical.}
\textit{Nature Astronomy} (2026).

\item Aparicio, V., Jones, D. O., Hoogendam, W. B., \textbf{Haghjoo, A.}, et al.
\textit{Analysis of a Type II Supernova Candidate at $z=3.19$ from JWST's COSMOS-Web Survey.}
\textit{Astrophysical Journal} (2026).

\item Abedini, F., Gozaliasl, G., Zonoozi, A. H., \textbf{Haghjoo, A.}, et al.
\textit{COSMOS-Web: Estimating Physical Parameters of Galaxies Using Self-Organizing Maps.}
\textit{Monthly Notices of the Royal Astronomical Society} (2026).

\item Hatamnia, H., Mobasher, B., Taamoli, S., \textbf{Haghjoo, A.}, et al.
\textit{Large-scale Structure in COSMOS-Web: Tracing Galaxy Evolution in the Cosmic Web up to $z\sim7$ with the Largest JWST Survey.}
\textit{Astrophysical Journal} (2026).

\item Shuntov, M., Ilbert, O., Lagos, C. del P., \textbf{Haghjoo, A.}, et al.
\textit{The Stellar Mass Function of Quiescent and Star-Forming Galaxies and its Dependence on Morphology in COSMOS-Web.}
\textit{Astronomy \& Astrophysics} (2026).

\item Yu, S.-Y., Ho, L. C., Tsukui, T., \textbf{Haghjoo, A.}, et al.
\textit{Through Thick and Thin: The Cosmic Evolution of Disk Scale Height.}
\textit{Astrophysical Journal Supplement Series} (2026).

\item Hemmati, S., Krick, J., Stern, D., \textbf{Haghjoo, A.}, et al.
\textit{Reducing the Dimensions of Active Galactic Nuclei Light-Curve Manifolds.}
\textit{Astrophysical Journal} (2026).

\item Toni, G., Maturi, M., Castignani, G., \textbf{Haghjoo, A.}, et al.
\textit{COSMOS-Web Galaxy Groups: Evolution of the Red Sequence and Quiescent Galaxy Fraction.}
\textit{Astronomy \& Astrophysics} (2026).

\item Fujimoto, S., Faisst, A. L., Tsujita, A., \textbf{Haghjoo, A.}, et al.
\textit{The ALPINE-CRISTAL-JWST Survey: NIRSpec IFU Data Processing and Spatially Resolved Views of Chemical Enrichment in Normal Galaxies at $z=4$--$6$.}
\textit{arXiv preprint} (2025).

\item Gozaliasl, G., Yang, L., Kartaltepe, J. S., \textbf{Haghjoo, A.}, et al.
\textit{Brightest Group Galaxies in COSMOS-Web: Evolution of the Size--Mass Relation Since $z=3.7$.}
\textit{Astronomy \& Astrophysics} (2025).

\item Asadi, V., Chartab, N., Zonoozi, A. H., \textbf{Haghjoo, A.}, et al.
\textit{Machine Learning Classification of COSMOS2020 Galaxies: Quiescent versus Star-Forming.}
\textit{Astrophysical Journal} (2025).

\end{itemize}

\end{section}
%--------------------------------------------------
\begin{section}{Research Experiences and Projects}

\begin{subsection}{Super-Resolving Low-Resolution Galaxy Spectra with Deep Learning}{Under the supervision of \href{https://faculty.ucr.edu/~mobasher/}{Bahram Mobasher} and \href{https://www.ipac.caltech.edu/meet-the-staff/shoubaneh-hemmati}{Shoubaneh Hemmati}}{Dec 2024 -- Present}{UCR}
    \item Developed a generative model to enhance spectral resolution of PRISM data by a factor of 10 using data from the JADES program
    \item Trained and validated the model with a physics-informed loss function targeting emission line accuracy and flux ratios (e.g., H$\alpha$/H$\beta$), achieving improved deblending of close spectral features
    \item Demonstrated the potential of the method to increase the resolution of deep spectral data from future instruments (e.g., Nancy Grace Roman Space Telescope)
    \item Currently extending the model to a multi-modal foundation framework that integrates photometric and spectroscopic data, enabling cross-modal inference
\end{subsection}

\begin{subsection}{Investigating the Effects of Non-Standard Physics on the Global 21cm Signal}{Under the supervision of \href{https://www.physics.mcgill.ca/~sievers/}{Jonathan Sievers} and \href{https://www.physics.mcgill.ca/~oscarh/}{Oscar Hernandez}}{Sept 2021 -- Dec 2023}{McGill University}
    \item Investigated the potential of the global 21cm signal to reveal non-standard physics
    \item Developed a pipeline using Markov Chain Monte Carlo (MCMC) and Levenberg-Marquardt (LM) algorithms to estimate the astrophysical parameters of the global signal
    \item Employed Accelerated Reionization Era Simulations (ARES) to generate theoretical models
    \item Validated the approach using data from Experiment to Detect the Global EoR Signature (EDGES)
    \item Results will aid in constraining future models and defining theoretical limits for upcoming experiments
\end{subsection}

% \begin{subsection}{A Machine Learning Approach to Classification of Variable Stars}{Under the supervision of \href{http://sharif.ir/~sraeisi/}{Sadegh Raeisi} and \href{https://sharif.edu/~rahvar/}{Sohrab Rahvar}}{Jan 2020 -- Sept 2020}{SUT}
%     \item Collaborated on the classification of variable stars in the OGLE dataset
%     \item Applied machine learning techniques including KNN, random forest, and CNNs
%     \item Addressed dataset imbalance challenges during model training
% \end{subsection}

%\begin{subsection}{Authoring ``Introduction to Cosmology" Book}{Under the supervision of \href{https://ipm.ac.ir/personalinfo.jsp?PeopleCode=IP1400002}{Bahram Mashhoon} \& \href{https://sharif.edu/~baghram/}{Shant Baghram}}{Jul 2020 - Jul 2021}{SUT}
 %   \item Collaborated on a project to author a public-audience cosmology book
  %  \item Transcribed one-third of Bahram Mashhoon's lectures on cosmology
  %  \item Replaced frequently used English cosmology terms with appropriate Persian equivalents
%\end{subsection}
\end{section}
%---------------------------------------------------
\sectiontable{Honours and Awards}{
    %\entry{Euclid Consortium Membership}{\hfill \textit{2025}}
    \entry{USC AI4SE Hackathon: Most Advanced Technology}{\hfill \textit{2026}}
    \entry{UCR Dean's Fellowship}{\hfill \textit{2024}}
    \entry{McGill Faculty of Science Graduate Award}{\hfill \textit{2023}}
    \entry{McGill Graduate Excellence Award}{\hfill \textit{2021}}
    \entry{Trottier Space Institute at McGill Graduate Fellowship}{\hfill \textit{2021}}
    %\entry{Lifetime Member of \href{https://en.bmn.ir/}{Iranian National Elites Foundation}}{\hfill \textit{2017}}
    \entry{National University Entrance Examinations: Top 0.5\% Among 180K Candidates}{\hfill \textit{2017}}
}

%--------------------------------------------------
\begin{section}{Presentations}

% Upcoming -- uncomment after Nov 12, 2026:
% \begin{subsectionnobullet}{From JWST to Roman: A Physics-Informed AI Super-Resolution Framework for Grism Spectroscopy}{Verbal presentation at the UCR Library AI Graduate Research Brown Bag Series}{Nov 2026}{UCR}
% \end{subsectionnobullet}

\begin{subsectionnobullet}{From JWST to Roman: A Physics-Informed AI Super-Resolution Framework for Grism Spectroscopy}{Verbal presentation at GalFRESCA 2026 (Galaxy Formation \& Evolution in Southern California)}{Sept 2026}{UC Irvine}
\end{subsectionnobullet}

\begin{subsectionnobullet}{From JWST to Roman: A Physics-Informed AI Super-Resolution Framework for Grism Spectroscopy}{Verbal presentation at \href{https://www.openskai-conference.org/}{Open SkAI 2026}, SkAI Institute}{Sept 2026}{Chicago, IL}
\end{subsectionnobullet}

\begin{subsectionnobullet}{From JWST to Roman: A Physics-Informed AI Super-Resolution Framework for Grism Spectroscopy}{\href{https://www.youtube.com/watch?v=z102gWq6Fo8\&t=2400s}{Verbal presentation} at the \href{https://www.stsci.edu/contents/events/stsci/2026/august/surveying-the-universe-in-4d-beating-cosmic-variance-with-wide-field-slitless-spectroscopy-from-hst}{Surveying the Universe in 4D Workshop} (virtual)}{Aug 2026}{Space Telescope Science Institute}
\end{subsectionnobullet}

\begin{subsectionnobullet}{From JWST to Roman: An AI Super-Resolution Framework for Grism Spectroscopy}{\href{https://www.youtube.com/watch?v=55EDsSILBbY}{Verbal presentation} at the \href{https://workshop.ipac.caltech.edu/roman2026/}{Shifting Landscapes in Astrophysics}}{Jul 2026}{IPAC, Caltech}
\end{subsectionnobullet}

\begin{subsectionnobullet}{Deep Learning for Astronomical Data Enhancement}{Verbal presentation at the \href{https://aas.org/meetings/aas248}{248th AAS Meeting}}{Jun 2026}{Pasadena, CA}
\end{subsectionnobullet}

\begin{subsectionnobullet}{Learning to See Sharper: A Physics-Informed AI Framework for Spectral Super-Resolution}{Poster presentation at the \href{https://datascience.stanford.edu/events/center-decoding-universe/c4du-annual-conference/2026-conference-physics-and-ai-pai26}{2026 Conference on Physics and AI (PAI26)}}{Jun 2026}{Stanford University}
\end{subsectionnobullet}

\begin{subsectionnobullet}{Learning to See Sharper: A Physics-Informed AI Framework for Spectral Super-Resolution}{\href{https://www.youtube.com/watch?v=M9KbXy-63ls\&t=1778s}{Verbal presentation} at the \href{https://astroai.cfa.harvard.edu/lunch_talks/}{AstroAI Lunch Talks} (Invited)}{Apr 2026}{CfA, Harvard and Smithsonian}
\end{subsectionnobullet}

\begin{subsectionnobullet}{Learning to See Sharper: A Physics-Informed AI Framework for Spectral Super-Resolution}{Verbal presentation at the Euclid US Tag-up Meeting (Invited)}{Mar 2026}{Euclid Consortium}
\end{subsectionnobullet}

\begin{subsectionnobullet}{Super-Resolving Low-Resolution Galaxy Spectra}{Poster presentation at the \href{https://tap.arizona.edu/news/arizona-winter-school-multi-probe-cosmology-roman-st-era}{Arizona Winter School on Cosmology in the Roman ST Era}}{Jan 2026}{University of Arizona}
\end{subsectionnobullet}

\begin{subsectionnobullet}{Super-Resolving Low-Resolution Galaxy Spectra}{Verbal presentation at the \href{https://conference.ipac.caltech.edu/IRSpaceSpectroscopy}{IR Spectroscopy from Space}}{Oct 2025}{IPAC, Caltech}
\end{subsectionnobullet}

\begin{subsectionnobullet}{Super-Resolving Low-Resolution Galaxy Spectra}{Verbal presentation at the \href{https://conference.ipac.caltech.edu/giss2025/}{Greater IPAC Science Symposium}}{Aug 2025}{IPAC, Caltech}
\end{subsectionnobullet}

\begin{subsectionnobullet}{Super-Resolving Low-Resolution Galaxy Spectra}{Poster presentation at the \href{https://www.stsci.edu/files/live/sites/www/files/home/events/event-assets/2025/_documents/2025-Roman-Symposium-Schedule.pdf}{Cosmic Cartography with Roman}}{Jul 2025}{Space Telescope Science Institute}
\end{subsectionnobullet}

\begin{subsectionnobullet}{Super-Resolving Low-Resolution Galaxy Spectra}{Verbal presentation at the \href{https://astroai.cfa.harvard.edu/workshop/details.html}{AstroAI Workshop}}{Jul 2025}{CfA, Harvard and Smithsonian}
\end{subsectionnobullet}

\begin{subsectionnobullet}{Unraveling Non-Standard Physics Through the Global 21cm Signal}{Verbal presentation at the \href{https://physics.sharif.edu/~cosmology/?p=2495}{SUT Cosmology Seminar}}{Jun 2023}{SUT}
\end{subsectionnobullet}

\begin{subsectionnobullet}{Parameter Estimation of the Global 21cm Signal}{\href{https://github.com/aryana-haghjoo/My-Project/blob/main/presentations/poster/poster.pdf}{Poster presentation} at \href{https://cosmosafari.co.za/}{Cosmology on Safari 2023}}{Mar 2023}{Hluhluwe, South Africa}
\end{subsectionnobullet}

\begin{subsectionnobullet}{Clarifying the Effects of Non-Standard Physics on the Global 21cm Signal}{Verbal presentation at the \href{https://www.physics.mcgill.ca/~chiang/}{McGill Radio Lab Meetings}}{Oct 2022}{McGill University}
\end{subsectionnobullet}

\begin{subsectionnobullet}{Effects of Non-Standard Physics on 21cm Signal}{Verbal presentation at the \href{https://www.physics.mcgill.ca/~rhb/}{McGill Early Universe Cosmology Meetings}}{May 2022}{McGill University}
\end{subsectionnobullet}

\end{section}


%--------------------------------------------------
\begin{section}{Work Experience}
\begin{subsectionnobullet}{Teaching Assistant}{PHYS 243 (Applied Machine Learning); PHYS 2C, 2B Discussions; PHYS 40A--C Labs}{2025 -- Present}{UCR}
    %\item PHYS 40C: General Physics (Apr 2025 -- Jun 2025)
    %\item PHYS 40B: General Physics (Jan 2025 -- Mar 2025)
\end{subsectionnobullet}

\begin{subsectionnobullet}{Teaching Assistant}{PHYS 102 Lab; Galaxies and Stars}{2022 -- 2023}{McGill University}
    %\item PHYS 557 (Sep 2023 -- Dec 2023)
    %\item PHYS 183 (Jan 2023 -- Apr 2023)
    %\item PHYS 102 (Jan 2022 -- Apr 2022)
\end{subsectionnobullet}

\begin{subsectionnobullet}{Teaching Assistant}{PHYS 103; Introduction to Cosmology; Electromagnetism}{2019 -- 2020}{SUT}
    %\item Introduction to Cosmology (Jan 2020 -- Apr 2020 \& Sept 2020 -- Dec 2020)
    %\item PHYS III (Sept 2019 -- Dec 2020)
    %\item Electromagnetism I (Jan 2020 -- Apr 2020)
\end{subsectionnobullet}
\end{section}
%--------------------------------------------------
\sectiontable{Skills and Languages}{
    \entry{Programming}{Python, MATLAB, \LaTeX, Bash, C, C++, Quick Basic, Visual Basic}
    \entry{Platforms}{GitHub, GitLab, SLURM}
    \entry{Languages}{Persian (native), English (fluent, full professional proficiency)}
}
%--------------------------------------------------
\sectiontable{Selected Courses}{
    \entry{Graduate}{Computational Physics, Galaxy Formation and Evolution, Astrophysics of the Interstellar Medium, Advanced Cosmology, Observational Astrophysics, Radiative Processes in Astrophysics, Stellar Structure and Evolution}
    \entry{Undergraduate}{Machine Learning in Physics, General Relativity, Quantum Field Theory, String Theory, Elementary Particles, Modeling Statistical Phenomena, Complex Systems}
}

%--------------------------------------------------

% \begin{section}{Conferences and Workshops}
% \begin{subsectionnobullet}{\href{https://exoplanetes.umontreal.ca/en/canadian-astronomical-instrumentation-workshop/}{Canadian Astronomical Instrumentation Workshop}}{In-person Participation}{Aug 2023}{\href{https://www.umontreal.ca/}{Université de Montréal}}
% \end{subsectionnobullet}

% \begin{subsectionnobullet}{\href{https://cosmosafari.co.za/}{Cosmology on Safari 2023}}{In-person Participation \& Presentation}{Mar 2023}{Hluhluwe, South Africa}
% \end{subsectionnobullet}

% \begin{subsectionnobullet}{\href{https://craq-astro.ca/summerschool/?lang=en}{CRAQ Summer School on Cosmology}}{In-person Participation}{Jun 2022}{McGill University}
% \end{subsectionnobullet}

% \begin{subsectionnobullet}{Introduction to Object-Oriented Coding for Physical Approaches Workshop}{Online Participation}{Mar 2020}{SUT}
% \end{subsectionnobullet}

% \begin{subsectionnobullet}{\href{http://www.psi.ir/farsi.asp?page=ngc98}{2020 Gravitation and Cosmology Conference of Iran}}{In-person Participation}{Jan 2020}{\href{https://ut.ac.ir/en}{University of Tehran}}
% \end{subsectionnobullet}

% \begin{subsectionnobullet}{\href{https://www.psi.ir/farsi.asp?page=physics98}{35th Annual Physics Conference of Iran}}{In-person Participation}{Aug 2019}{\href{https://tabrizu.ac.ir/en}{University of Tabriz}}
% \end{subsectionnobullet}

% \begin{subsectionnobullet}{Data Science and Data Modeling in Physics Workshop}{In-person Participation}{Aug 2019}{\href{http://tabrizu.ac.ir/en}{University of Tabriz}}
% \end{subsectionnobullet}

% \begin{subsectionnobullet}{Machine Learning in Physics Workshop}{In-person Participation}{Aug 2018}{\href{http://ikiu.ac.ir/en/}{Imam Khomeini International University}}
% \end{subsectionnobullet}
% \end{section}
% \vspace{0.5cm}

%--------------------------------------------------
\sectiontable{Other Activities and Interests}{
    \entry{Mentor at STEM High School: Assisting Senior Students on Their Projects}{\hfill \textit{Sept 2025 -- Mar 2026}}
    \entry{Member, \href{https://www.ted.com/tedx/events/53037}{TEDxAveLorne} Curation Team}{\hfill \textit{Jul 2023 -- Dec 2023}}
    \entry{Representative, \href{https://mgaps.physics.mcgill.ca/}{McGill Graduate Association of Physics Students (MGAPS)}}{\hfill \textit{Jan 2023 -- Dec 2023}}
    \entry{SUT Mentoring Program Volunteer: Senior-Freshman Buddy Program}{\hfill \textit{2020 -- 2021}}
    \entry{Amateur Astronomer: Outreach Observing Nights}{\hfill \textit{2015 -- 2020}}
}
%--------------------------------------------------
\sectiontable{References}{
    \entry{Bahram Mobasher}{Distinguished Professor, UCR \hfill \email{mobasher@ucr.edu}}
    \entry{Shoubaneh Hemmati}{Research Scientist, Caltech \hfill \email{shemmati@ipac.caltech.edu}}
    \entry{Jonathan Sievers}{Associate Professor, McGill University \hfill \email{jonathan.sievers@mcgill.ca}}
    \entry{Oscar Hernandez}{Adjunct Professor, McGill University \hfill \email{oscarh@physics.mcgill.ca}}
    \entry{Bahram Mashhoon}{Emeritus Professor, University of Missouri \hfill \email{mashhoonb@missouri.edu}}
    \entry{Shant Baghram}{Associate Professor, SUT \hfill \email{baghram@sharif.edu}}
}

\end{document}
